% Main-text section: expectation values.  Input file; no preamble.  Labels prefix ev.
\section{Expectation values}\label{sec:expect}

The theorems compare the solution with the periodized approximation in $L^2(M)$. This section records what follows for expectation
values, and shows on the modular surface that accuracy of the approximation and classical behaviour of expectation values are
different things: the approximation stays accurate up to times of order $\hbar^{-1/2}$, while the classical transport of the
initial density fails to describe the expectation values after a time of order $\log(1/\hbar)\log\log(1/\hbar)$.

\subsection{Transfer and cancellation}\label{subsec:ev-transfer}
\begin{corollary}\label{ev:transfer}
Assume the hypotheses of Theorem~\ref{thm:dd14}, or those of Theorem~\ref{pc:main} with $\pcNorm a$ in place of $\norm a$. Let
$K\in\mathbb N_0$ and $0\le t\le\delta(\alpha/\hbar)^{1/2}$, and put
\[
  \psi(t):=\cU(t)\big[a\,\ue^{i\varphi_b/\hbar}\big]_\Gamma,\qquad
  \epsilon_K:=C^{K+1}K!\,\Big(\frac{\hbar t}{\alpha^2}\Big)^{K+1}\norm a .
\]
\begin{itemize}
\item[(i)] For every bounded operator $B$ on $L^2(M)$,
\[
  \big|\langle B\psi(t),\psi(t)\rangle-\langle Bu^{(K)}(t)_\Gamma,u^{(K)}(t)_\Gamma\rangle\big|
  \le\norm B\,\epsilon_K\big(2\norm{\psi(0)}+\epsilon_K\big).
\]
\item[(ii)] $\big|\,\norm{u^{(K)}(t)_\Gamma}_{L^2(M)}-\norm{[a\,\ue^{i\varphi_b/\hbar}]_\Gamma}_{L^2(M)}\big|\le\epsilon_K$.
\end{itemize}
\end{corollary}
\begin{proof}
The theorem at $N=0$ gives $\norm{\psi(t)-u^{(K)}(t)_\Gamma}\le\epsilon_K$, and $\norm{\psi(t)}=\norm{\psi(0)}$ since $\cU(t)$ is unitary.
This is (ii). For (i) write $u=u^{(K)}(t)_\Gamma$ and
$\langle B\psi,\psi\rangle-\langle Bu,u\rangle=\langle B(\psi-u),\psi\rangle+\langle Bu,\psi-u\rangle$, and use
$\norm u\le\norm{\psi(0)}+\epsilon_K$.
\end{proof}

\begin{remark}[cancellation]\label{ev:cancel}
At $K=0$, (ii) says that $\norm{[(S(t)a)\,\ue^{i\varphi_b/\hbar}]_\Gamma}_{L^2(M)}$ stays within $C\hbar t\alpha^{-2}\norm a$ of its value at $t=0$,
for all $t\le\delta(\alpha/\hbar)^{1/2}$. Without the phase this fails: if $M$ has finite area and $a\in L^1(\Hb)$, then by the
Cauchy--Schwarz inequality and unfolding
\[
  \norm{[S(t)|a|]_\Gamma}_{L^2(M)}\ge\operatorname{area}(M)^{-1/2}\int_\Hb S(t)|a|\,d\nu=\operatorname{area}(M)^{-1/2}\,\ue^{t/2}\norm a_{L^1(\Hb)} .
\]
So the norm of the periodization with the phase is smaller than the norm of the periodization of the moduli by a factor of order
at most $\ue^{-t/2}$: the terms cancel. The following picture is geometric intuition and is not proved here. The terms that
contribute at a point come from about $\ue^{t}$ sheets of the cover. Up to the Ehrenfest time $\log(1/\hbar)$ their phase
gradients differ by more than $\hbar$, so that they are almost orthogonal one by one; beyond it the phase gradients of many pairs
differ by less than $\hbar$, and the cancellation is a property of the sum.
\end{remark}

By (i), for bounded observables the expectation values of $\psi(t)$ up to times $\delta(\alpha/\hbar)^{1/2}$ are those of the
explicit sum $u^{(K)}(t)_\Gamma$, up to an error that tends to $0$ with $\hbar$ for fixed $\alpha$ and $K$. The next subsection answers one such question in a case where the sum can be evaluated.

\subsection{The modular surface}\label{subsec:ev-mod}
In this subsection $\Gamma=\mathrm{PSL}(2,\mathbb Z)$, $M=\Gamma\backslash\Hb$, $x_0=i$, $b=\infty$ and $\varphi=-\log y$. Let
$F=\{|x|\le\tfrac12,\ |z|\ge1\}$ be the standard fundamental domain; for $Y\ge1$ let $M_Y\subset M$ be the image of $\{z\in F:y\le Y\}$,
a compact set, and $C_Y$ the image of $\{z\in F:y>Y\}=\{|x|\le\tfrac12,\ y>Y\}$, the part of the cusp above height $Y$. Recall
$\operatorname{area}(M)=\pi/3$.

Fix $\alpha\in(0,\tfrac12]$ and $h\in C_c^\infty(\Hb)$, $h\ge0$, such that
\[
  h_0(y):=\sum_{n\in\mathbb Z}h(x+n,y)
\]
does not depend on $x$ and is not identically zero; for instance $h(x,y)=\int_0^1h_1(x+s,y)\,ds$ with $h_1\in C_c^\infty(\Hb)$,
$h_1\ge0$, $h_1\not\equiv0$. Let
\[
  a:=\ue^{-\alpha\Lambda}h .
\]
Then $\pcNorm a=\norm{\ue^{\beta\langle d\rangle}h}<\infty$ for every $\beta$, so $a$ is an amplitude of the printed class, and it does
not depend on $\hbar$. Let $k_\alpha\ge0$ be the radial kernel of $\ue^{-\alpha\Lambda}$ (Lemma~\ref{os:poisson}). It is integrable on
$\Hb$, by \textup{(\ref{os:ub})} and \textup{(\ref{os:small})}. Since $k_\alpha(d(z,w))$ is invariant under $z\mapsto z+n$, $w\mapsto w+n$,
\begin{equation}\label{ev:a0}
  a_0(y):=\sum_{n\in\mathbb Z}a(x+n,y)=\int_\Hb k_\alpha(d(z,w))\,h_0(\operatorname{Im}w)\,d\nu(w)
\end{equation}
does not depend on $x$, and $a_0>0$. Put
\[
  \mu_0:=\int_0^\infty|a_0(y)|^2\,\frac{dy}{y^2}.
\]

\begin{proposition}[accurate approximation, non-classical expectation values]\label{ev:mod}
Let $\Gamma$, $\alpha$, $h$, $a$, $a_0$, $\mu_0$ be as above, $\psi(t)=\cU(t)[a\,\ue^{i\varphi/\hbar}]_\Gamma$, and let $C,\delta$ be the
constants of Theorem~\ref{pc:main} for $\beta=1$ and $\kappa=1$. Then $0<\mu_0<\infty$, and there is an absolute constant $C_1$ such
that, with $T(\hbar):=C_1\log(1/\hbar)\log\log(1/\hbar)$, the following hold.
\begin{itemize}
\item[(i)] $\norm{\psi(t)}^2_{L^2(M)}=\mu_0+O(\hbar^{n})$ for every $n$, and for all $K$ and $0\le t\le\delta(\alpha/\hbar)^{1/2}$,
$\norm{\psi(t)-u^{(K)}(t)_\Gamma}_{L^2(M)}\le C^{K+1}K!\,(\hbar t/\alpha^2)^{K+1}\pcNorm a$.
\item[(ii)] For every $Y\ge1$ and every $K$, as $\hbar\to0$,
\[
  \int_{M_Y}|\psi(t)|^2\,d\nu+\int_{M_Y}|u^{(K)}(t)_\Gamma|^2\,d\nu\longrightarrow0
  \qquad\text{uniformly for }T(\hbar)+\log Y\le t\le\delta(\alpha/\hbar)^{1/2}.
\]
\item[(iii)] For every continuous $q$ with compact support in $M$, lifted to a $\Gamma$-invariant function on $\Hb$,
\[
  \int_0^\infty\!\!\int_0^1q(z)\,\big|(S(t)a_0)(y)\big|^2\,dx\,\frac{dy}{y^2}\longrightarrow\frac3\pi\int_Mq\,d\nu\cdot\mu_0\qquad(t\to\infty).
\]
\end{itemize}
\end{proposition}
The left side of (iii) is the classical prediction for $\langle q\,\psi(t),\psi(t)\rangle$: the density $|a_0|^2$ on
$\Gamma_\infty\backslash\Hb$, transported by the geodesic flow for the time $t$ and projected to $M$; it is the sum of the squared
moduli of the terms of $u^{(0)}(t)_\Gamma$, summed over $\Gamma_\infty\backslash\Gamma$. By (iii) this prediction tends to the
equidistributed value. By (ii), for $q\ge0$ with support in $M_Y$,
$\langle q\,\psi(t),\psi(t)\rangle\le\sup q\int_{M_Y}|\psi(t)|^2\to0$ uniformly on the window of (ii): the $L^2$ mass of the state in every
compact set tends to $0$; the state has moved into the cusp. So on this window the expectation value of a compactly supported
$q\ge0$, $q\not\equiv0$, is not given by the classical transport, whose limit $\tfrac3\pi\mu_0\int_Mq\,d\nu$ is positive, although by (i) and
(ii) the periodized approximation remains accurate and has the same behaviour. The time $T(\hbar)$ is a sufficient time; it is not
claimed to be the time at which the departure begins. For the wave
group and Lagrangian states on unstable horocycles of the modular surface, the departure from classical transport at logarithmic
times is due to Brooks \cite{Bro17}.

\begin{remark}[the time scale]\label{ev:time}
The proof uses the classical zero-free region of the Riemann zeta function, which gives the sufficient time $T(\hbar)$. For $r>0$
and a continuous choice of the argument, the scattering coefficient $\Phi$ of \eqref{ev:Phi} satisfies
$-\tfrac{d}{dr}\arg\Phi(\tfrac12+ir)=2\log(r/\pi)+4\operatorname{Re}\tfrac{\zeta'}{\zeta}(1+2ir)+O(r^{-2})$. Heuristically this is the time delay of a wave
packet at $r=1/\hbar$, and its first term is twice the Ehrenfest time $\log(1/\hbar)$, up to a constant. The bound for $\zeta'/\zeta$ used here is
of the same order $\log r$ as that term, so this is an indication only: we do not prove a statement with the constant $2$, and we
do not prove that the classical prediction holds before that time.
\end{remark}

\subsubsection*{Proof of Proposition~\ref{ev:mod}}
Throughout, $n$ denotes an arbitrary positive integer and constants may depend on $n$, $\alpha$ and $h$.

\emph{Step 1: the function $a_0$.} Let $w$ be the bounded function $w(z)=h_0(\operatorname{Im}z)$ on $\Hb$. Since $k_\alpha\ge0$ is
integrable, \eqref{ev:a0} gives $0<a_0\le\norm{k_\alpha}_{L^1(\Hb)}\sup h_0$, and by Tonelli's theorem
\[
  \int_0^\infty a_0(y)\frac{dy}{y^2}=\int_0^\infty\!\!\int_0^1\sum_na(x+n,y)\,dx\frac{dy}{y^2}=\int_\Hb a\,d\nu\le\norm{k_\alpha}_{L^1}\norm h_{L^1}<\infty .
\]
Hence $\mu_0\le\sup a_0\int_0^\infty a_0\,dy/y^2<\infty$, and $\mu_0>0$. Convolution with a radial kernel commutes with the isometries
$z\mapsto\lambda z$, $\lambda>0$, whose generator acts on functions of $y$ alone as $y\partial_y$. So
$(y\partial_y)^ja_0$ is given by \eqref{ev:a0} with $(y\partial_y)^jh_0$ in place of $h_0$, and since
$|(y\partial_y)^jh_0|\le C_j\tilde h_0$ for a fixed $\tilde h_0\in C_c^\infty((0,\infty))$, $\tilde h_0\ge0$, that equals $1$ on the support of $h_0$, the same
bounds hold: with $\tilde a_0$ defined by \eqref{ev:a0} from $\tilde h_0$,
\begin{equation}\label{ev:a0bd}
  |(y\partial_y)^ja_0|\le C_j\tilde a_0,\qquad \sup\tilde a_0<\infty,\qquad\int_0^\infty\tilde a_0\,\frac{dy}{y^2}<\infty\qquad(j\ge0).
\end{equation}
For a function $f$ on $(0,\infty)$ let $\hat f(s):=\int_0^\infty f(y)y^{-s}\,dy/y$ be its Mellin transform. By \eqref{ev:a0bd},
$\hat a_0(s)$ converges absolutely for $0<\operatorname{Re}s\le1$, is holomorphic in $0<\operatorname{Re}s<1$ and continuous up to
$\operatorname{Re}s=1$. In the variable $u=\log y$ the functions $y^{-\sigma}(y\partial_y)^ja_0$, $\tfrac14\le\sigma\le1$, are integrable with integrable
derivatives, so they vanish at $\pm\infty$, and integration by parts gives $s^j\hat a_0(s)=\widehat{(y\partial_y)^ja_0}(s)$. Hence
\begin{equation}\label{ev:mellin}
  |\hat a_0(\sigma+i\rho)|\le C_n(1+|\rho|)^{-n}\qquad(\tfrac14\le\sigma\le1,\ \rho\in\mathbb R).
\end{equation}
The Mellin--Plancherel formula reads $\int_0^\infty|f|^2\,dy/y^2=\tfrac1{2\pi}\int_{\mathbb R}|\hat f(\tfrac12+ir)|^2\,dr$.

\emph{Step 2: incomplete Eisenstein series.} For a function $f$ on $(0,\infty)$ put $E(z|f):=\sum_{\gamma\in\Gamma_\infty\backslash\Gamma}f(\operatorname{Im}\gamma z)$,
where $\Gamma_\infty$ is the group of the translations $z\mapsto z+n$. The cosets are the identity and the classes of the
matrices with bottom row $(c,d)$, $c\ge1$, $\gcd(c,d)=1$, and $\operatorname{Im}\gamma z=y/|cz+d|^2$. Let $\mathcal F$ be the set of smooth $f$ with
$\int_0^\infty|f|\,dy/y^2<\infty$ and $\int_0^\infty|f|^2\,dy/y^2<\infty$. For $f\in\mathcal F$, unfolding gives
$\int_ME(z||f|)\,d\nu=\int_0^\infty|f|\,dy/y^2<\infty$; so the series converges absolutely almost everywhere and $E(\cdot|f)\in L^1(M)$. For $y_1>0$ the
set $\{0\le x\le1,\ y>y_1\}$ is covered by finitely many, say $\mathcal N(y_1)$, $\Gamma$-translates of $F$, with $\mathcal N(1)=2$. So
$\int_{y_1}^\infty\int_0^1E(z||f|)\,d\nu\le\mathcal N(y_1)\int_0^\infty|f|\,dy/y^2$, and for almost every $y>0$ the zero Fourier mode of $E(\cdot|f)$ is, by
Fubini's theorem and the substitution $d=d_0+cn$,
\begin{equation}\label{ev:zero}
  \int_0^1E(x+iy|f)\,dx=f(y)+R_f(y),\qquad R_f(y):=\sum_{c\ge1}N_c\int_{\mathbb R}f\Big(\frac{y}{c^2(x^2+y^2)}\Big)dx ,
\end{equation}
where $N_c$ is the number of residues $d_0$ modulo $c$ that are prime to $c$. The same computation with $|f|$ shows
$\int_{y_1}^\infty|R_f|\,dy/y^2\le\mathcal N(y_1)\int_0^\infty|f|\,dy/y^2$ for every $y_1>0$. Let
\begin{equation}\label{ev:Phi}
  \Phi(s):=\sqrt\pi\,\frac{\Gamma(s-\frac12)}{\Gamma(s)}\,\frac{\zeta(2s-1)}{\zeta(2s)}
  =\pi^{2s-1}\,\frac{\Gamma(1-s)}{\Gamma(s)}\,\frac{\zeta(2-2s)}{\zeta(2s)} ,
\end{equation}
the two forms being equal by the functional equation of $\zeta$. $\Phi$ is holomorphic in $\operatorname{Re}s\ge\tfrac12$ except for a simple
pole at $s=1$ with residue $3/\pi$, $|\Phi(\tfrac12+ir)|=1$ (because $\Phi(s)\Phi(1-s)=1$ and $\overline{\Phi(s)}=\Phi(\bar s)$), and $\Phi$ is bounded by a power of $1+|s|$ in
$\tfrac12\le\operatorname{Re}s\le2$, $|s-1|\ge\tfrac14$: the Gamma quotient by Stirling's formula \cite[5.11.9]{DLMF}, $\zeta(2s-1)$ by its polynomial
growth in vertical strips, and $1/\zeta(2s)$ by \cite[Theorem~6.7]{MV07}.
We claim that for $f\in\mathcal F$, for almost every $y>0$,
\begin{equation}\label{ev:R}
  R_f(y)=\frac3\pi\hat f(1)+\mathcal O_f(y),\qquad \mathcal O_f(y):=\frac1{2\pi}\int_{\mathbb R}\Phi(\tfrac12+ir)\hat f(\tfrac12+ir)\,y^{1/2-ir}\,dr ,
\end{equation}
where the integral is the $L^2$ inverse Mellin transform, so that
\begin{equation}\label{ev:Onorm}
  \int_0^\infty|\mathcal O_f|^2\,\frac{dy}{y^2}=\int_0^\infty|f|^2\,\frac{dy}{y^2}.
\end{equation}
For $f\in C_c^\infty((0,\infty))$, $\hat f$ is entire and rapidly decreasing on vertical lines, and Mellin inversion on $\operatorname{Re}s=2$ gives
\[
  R_f(y)=\frac1{2\pi i}\int_{(2)}\hat f(s)\Big(\sum_{c\ge1}\frac{N_c}{c^{2s}}\Big)\Big(\int_{\mathbb R}\frac{y^s\,dx}{(x^2+y^2)^s}\Big)ds
  =\frac1{2\pi i}\int_{(2)}\hat f(s)\,\Phi(s)\,y^{1-s}\,ds ,
\]
because $\sum_cN_cc^{-2s}=\zeta(2s-1)/\zeta(2s)$ and $\int_{\mathbb R}(1+v^2)^{-s}dv=\sqrt\pi\,\Gamma(s-\tfrac12)/\Gamma(s)$. Moving the line to
$\operatorname{Re}s=\tfrac12$ gives \eqref{ev:R}. For general $f\in\mathcal F$ take $f_j=f\chi(j^{-1}\log y)$ with $\chi\in C_c^\infty(\mathbb R)$, $0\le\chi\le1$, $\chi=1$ near $0$:
then $f_j\to f$ in $L^1$ and in $L^2$ of $dy/y^2$, so $R_{f_j}\to R_f$ in $L^1((y_1,\infty),dy/y^2)$ for every $y_1>0$, $\hat f_j(1)\to\hat f(1)$ and
$\mathcal O_{f_j}\to\mathcal O_f$ in $L^2(dy/y^2)$, and \eqref{ev:R} follows.

\emph{Step 3: the approximation as an incomplete Eisenstein series.} Let
\[
  f_t(y):=(S(t)a_0)(y)\,y^{-i/\hbar}=\ue^{-t/2}a_0(\ue^ty)\,y^{-i/\hbar}.
\]
By \eqref{ev:a0bd}, $f_t\in\mathcal F$, with $\int_0^\infty|f_t|^2dy/y^2=\mu_0$ and $\int_0^\infty|f_t|\,dy/y^2=\ue^{t/2}\int_0^\infty a_0\,dy/y^2$. By Corollary~\ref{os:cwp}, $\omega a=\omega\,\ue^{-\alpha\Lambda}h\in L^2(\Hb)$ for the periodization weight $\omega$ of
Section~\ref{sec:printed}. So by Proposition~\ref{os:p1} and (W1), (W2) of Section~\ref{subsec:dd-surf}, the orbit sums of $S(t)a$ and of
$(S(t)a)\ue^{i\varphi/\hbar}$ converge absolutely almost everywhere and lie in $L^2(M)$. They may be summed over $\Gamma_\infty$ first, so
$E(\cdot||f_t|)=[S(t)a]_\Gamma\in L^2(M)$ and
\[
  \big[(S(t)a)\ue^{i\varphi/\hbar}\big]_\Gamma=E(\cdot|f_t),\qquad u^{(0)}(t)_\Gamma=\ue^{-it/(2\hbar)}E(\cdot|f_t),\qquad\psi(0)=E(\cdot|f_0).
\]
The substitution $\eta=\ue^ty$ gives
\begin{equation}\label{ev:ft}
  \hat f_t(s)=\ue^{t(s-1/2)}\,\ue^{it/\hbar}\,\hat a_0(s+i/\hbar)\qquad(0<\operatorname{Re}s\le1).
\end{equation}
Theorem~\ref{pc:main} with $\beta=\kappa=1$ applies for $\hbar\le\alpha^3$; write $\epsilon_0=C\hbar t\alpha^{-2}\pcNorm a$. For
$t\le\delta(\alpha/\hbar)^{1/2}$ we have $\epsilon_0\le C\delta\alpha^{-3/2}\hbar^{1/2}\pcNorm a\to0$, and the second statement of (i) is the theorem.

\emph{Step 4: the norm and the mean.} Since $E(\cdot||f_0|)\in L^2(M)$, unfolding one factor of $|E(\cdot|f_0)|^2$ is justified, and with
\eqref{ev:zero} and \eqref{ev:R} it gives
\[
  \norm{\psi(0)}^2=\int_0^\infty f_0\,\overline{\Big(f_0+\frac3\pi\hat f_0(1)+\mathcal O_{f_0}\Big)}\,\frac{dy}{y^2}
  =\mu_0+\frac3\pi|\hat f_0(1)|^2+\int_0^\infty f_0\,\overline{\mathcal O_{f_0}}\,\frac{dy}{y^2},
\]
where the three integrals converge absolutely because $f_0\in\mathcal F$ and $\mathcal O_{f_0}\in L^2(dy/y^2)$, and $\int_0^\infty f_0\,dy/y^2=\hat f_0(1)$.
By \eqref{ev:ft} and \eqref{ev:mellin}, $|\hat f_0(1)|=|\hat a_0(1+i/\hbar)|\le C_n\hbar^n$. The Mellin transform of $\mathcal O_{f_0}$ at
$\tfrac12+ir$ is $\Phi(\tfrac12-ir)\hat f_0(\tfrac12-ir)$, so by the Parseval formula and \eqref{ev:ft} the last integral is at most
$\tfrac1{2\pi}\int|\hat a_0(\tfrac12+i(r+\tfrac1\hbar))|\,|\hat a_0(\tfrac12+i(-r+\tfrac1\hbar))|\,dr\le C_n\hbar^n$: at every $r$ one of the two arguments has
imaginary part at least $1/\hbar$ in modulus, and \eqref{ev:mellin} is used with the power $n+2$ for that factor and with the power $2$ for the
other, which is integrable in $r$. Hence $\norm{\psi(t)}^2=\norm{\psi(0)}^2=\mu_0+O(\hbar^n)$, which is the first statement of (i).

Since $M$ has finite area, the constant function $1$ lies in $L^2(M)$ and $\cU(t)1=1$. So $\langle\psi(t),1\rangle=\langle\psi(0),1\rangle=\hat f_0(1)$ and
$\langle E(\cdot|f_t),1\rangle=\hat f_t(1)$ (unfolding), and by the theorem
\begin{equation}\label{ev:mean}
  |\hat f_t(1)|\le|\hat f_0(1)|+\operatorname{area}(M)^{1/2}\epsilon_0\le C_n\hbar^n+2\epsilon_0\qquad(0\le t\le\delta(\alpha/\hbar)^{1/2}).
\end{equation}

\emph{Step 5: the outgoing term below height $Y$.} We bound $\int_0^Y|\mathcal O_{f_t}|^2dy/y^2$. Let $c_0\in(0,1]$ be strictly smaller than the constant $c$ of \cite[Theorem~6.6]{MV07}. Then the closed region
$\operatorname{Re}w\ge1-c_0/(2\log(|\operatorname{Im}w|+4))$ has a neighbourhood in which $\zeta(w)\ne0$, and in this region
$|\log\zeta(w)|\le\log\log(|\operatorname{Im}w|+4)+C$ where $|\operatorname{Im}w|\ge\tfrac78$, and $\zeta(w)(w-1)$ is bounded above and away from $0$ where
$|\operatorname{Im}w|\le\tfrac78$ \cite[Theorems~6.6 and~6.7]{MV07}. Put
\[
  \vartheta(r):=\frac{c_0}{4\log(2|r|+4)}\in(0,\tfrac14),\qquad D:=\{s=\sigma+ir:\ \tfrac12-\vartheta(r)\le\sigma\le\tfrac12\}.
\]
For $s\in D$ the points $2s$ and $2-2s$ lie in that region. So for $|r|\ge1$, $|\zeta(2-2s)|$ and $|1/\zeta(2s)|$ are at most
$C\log(|r|+3)$; and $|\Gamma(1-s)/\Gamma(s)|\le C(1+|r|)^{1-2\sigma}\le C(1+|r|)^{2\vartheta(r)}\le C'$ by Stirling's formula \cite[5.11.9]{DLMF}. For
$|r|\le1$, $\Phi$ is holomorphic on a neighbourhood of the compact set $D\cap\{|r|\le1\}$: $\zeta(2s)$ has no zero there, and in the
quotient $\zeta(2-2s)/\zeta(2s)$ the two poles at $s=\tfrac12$ cancel. So by the second form in \eqref{ev:Phi}, $\Phi$ is holomorphic on a neighbourhood of $D$ and
\begin{equation}\label{ev:PhiD}
  |\Phi(s)|\le C\log^2(|r|+3)\qquad(s\in D).
\end{equation}
By \eqref{ev:ft} and \eqref{ev:mellin}, on $D$ the function $\hat f_t$ is holomorphic and
$|\hat f_t(\sigma+ir)|\le\ue^{-t(1/2-\sigma)}C_n(1+|r+\tfrac1\hbar|)^{-n}$. So in the integral that defines $\mathcal O_{f_t}(y)$, written as
$\tfrac1{2\pi i}\int_{(1/2)}\Phi(s)\hat f_t(s)y^{1-s}ds$, the line may be moved to the left edge $\sigma=\tfrac12-\vartheta(r)$ of $D$, a piecewise smooth curve: for $f_t$ the integral converges
absolutely and defines the $L^2$ inverse Mellin transform pointwise, and the integrals over the horizontal segments $\{\sigma+ir:\tfrac12-\vartheta(r)\le\sigma\le\tfrac12\}$
tend to $0$ as $|r|\to\infty$, by \eqref{ev:PhiD} and the decay of $\hat f_t$. On that edge $|y^{1-s}|=y^{1/2+\vartheta(r)}$, the
norm of $y^{1/2+\vartheta}$ in $L^2((0,Y),dy/y^2)$ is $Y^{\vartheta}(2\vartheta)^{-1/2}$, and $|ds|\le2\,dr$. Minkowski's inequality gives, for $t\ge\log Y$,
\[
  \Big(\int_0^Y|\mathcal O_{f_t}|^2\frac{dy}{y^2}\Big)^{1/2}
  \le C_n\int_{\mathbb R}\log^{5/2}(|r|+4)\,\ue^{-(t-\log Y)\vartheta(r)}\,(1+|r+\tfrac1\hbar|)^{-n}\,dr .
\]
On $|r+\tfrac1\hbar|\le\tfrac1{2\hbar}$ we have $\vartheta(r)\ge\vartheta_\hbar:=c_0/(4\log(3/\hbar+4))$; on the rest the exponential is at most $1$ and
the integral is $O(\hbar^{n-2})$. So
\begin{equation}\label{ev:out}
  \Big(\int_0^Y|\mathcal O_{f_t}|^2\frac{dy}{y^2}\Big)^{1/2}\le C\log^{5/2}(4/\hbar)\,\ue^{-(t-\log Y)\vartheta_\hbar}+C_n\hbar^{n-2}.
\end{equation}
With $C_1:=13/c_0$ and $t-\log Y\ge T(\hbar)$, the exponent is at least $\tfrac{13}{4}\log\log(1/\hbar)\cdot\log(1/\hbar)/\log(3/\hbar+4)$, and the right side of
\eqref{ev:out} tends to $0$ as $\hbar\to0$, uniformly in such $t$ and in $Y\ge1$.

\emph{Step 6: the mass in the cusp.} Let $Y\ge1$ and $T(\hbar)+\log Y\le t\le\delta(\alpha/\hbar)^{1/2}$, and write $v=E(\cdot|f_t)$. By the Cauchy--Schwarz
inequality in $x$, \eqref{ev:zero} and \eqref{ev:R},
\begin{align*}
  \Big(\int_{C_Y}|v|^2d\nu\Big)^{1/2}&\ge\Big(\int_Y^\infty\Big|\int_0^1v\,dx\Big|^2\frac{dy}{y^2}\Big)^{1/2}\\
  &\ge\norm{\mathcal O_{f_t}}_{L^2((Y,\infty))}-\norm{f_t}_{L^2((Y,\infty))}-\frac3\pi|\hat f_t(1)|\,Y^{-1/2},
\end{align*}
with the norms of $L^2(dy/y^2)$. By \eqref{ev:Onorm} and \eqref{ev:out}, $\norm{\mathcal O_{f_t}}^2_{L^2((Y,\infty))}=\mu_0-o(1)$. Next
$\norm{f_t}^2_{L^2((Y,\infty))}=\int_{\ue^tY}^\infty|a_0|^2d\eta/\eta^2\le\sup|a_0|^2\,\ue^{-t}\to0$, and $|\hat f_t(1)|\to0$ by \eqref{ev:mean}. So
$\int_{C_Y}|v|^2d\nu\ge\mu_0-o(1)$. By Corollary~\ref{ev:transfer}(ii) and Step 4, $\norm v^2_{L^2(M)}\le(\norm{\psi(0)}+\epsilon_0)^2=\mu_0+o(1)$. Hence
\[
  \int_{M_Y}|u^{(0)}(t)_\Gamma|^2d\nu=\norm v^2-\int_{C_Y}|v|^2d\nu=o(1),
\]
uniformly on the window. Then $\norm{\mathbf 1_{M_Y}\psi(t)}\le\norm{\mathbf 1_{M_Y}u^{(0)}(t)_\Gamma}+\epsilon_0$ and
$\norm{\mathbf 1_{M_Y}u^{(K)}(t)_\Gamma}\le\norm{\mathbf 1_{M_Y}\psi(t)}+\epsilon_K$, with $\epsilon_K=C^{K+1}K!(\hbar t/\alpha^2)^{K+1}\pcNorm a\to0$ on the window for
fixed $K$. This proves (ii).

\emph{Step 7: the classical transport.} For $q$ continuous with compact support in $M$, the closed horocycles equidistribute:
$\int_0^1q(x+iy)\,dx\to\tfrac3\pi\int_Mq\,d\nu$ as $y\to0$. For smooth $q$ with compact support this is in \cite[Section~1, pp.~279--280]{Zag81}
(see also \cite{Sar81}); for continuous $q$ it follows by uniform approximation, since the measures $dx$ on the horocycles and
$\tfrac3\pi d\nu$ have mass $1$. The left side of (iii) equals
$\int_0^\infty|a_0(\eta)|^2\big(\int_0^1q(x+i\ue^{-t}\eta)\,dx\big)\,d\eta/\eta^2$, and (iii) follows by dominated convergence. \qed
